\documentclass[a4paper]{article}
\usepackage{MMVII}

\newcommand{\TheCmd}{EditProfile}

\begin{document}
\sloppy

\CmdTitle

\tableofcontents


 %=================================================================================

\SecDescr

This command inspects and edits MMVII user profiles.  A user profile is a set of
key/value preferences used by MMVII commands, for example the default serialization
formats, the maximum number of parallel processes, or the command used to open PDF
files.

When \TheCmd\ is called without option, it prints:

\begin{itemize}
  \item the known profile names, with the current profile shown between square brackets;
  \item the path of the profile file effectively loaded;
  \item the key/value pairs of the current profile.
\end{itemize}

\begin{verbatim}
MMVII EditProfile
\end{verbatim}


 %=================================================================================

\section{Selecting the current profile}

The optional parameter \FileSty{SetCurrent} makes a profile the default profile for
subsequent MMVII commands:

\begin{verbatim}
MMVII EditProfile SetCurrent=Default
MMVII EditProfile SetCurrent=Test
\end{verbatim}

The profile name must contain only letters, digits, underscores, or hyphens.  It
must not be empty.  If the selected profile file does not exist yet, MMVII creates
it by synchronizing it with the distributed default profile.

The current profile name is stored in the user configuration file
\FileSty{MMVII-Current-Profile.xml}.


 %=================================================================================

\section{Editing profile values}

The optional parameter \FileSty{KeyVal} sets one key/value pair in the selected
profile.  It must contain exactly two values: the key and the value.
To specify an empty string for a value, use the special syntax  \FileSty{KeyVal=[Key,,]}.

\begin{verbatim}
MMVII EditProfile KeyVal=[NbProcMax,8]
MMVII EditProfile KeyVal=[TaggedSerialMode,json]
MMVII EditProfile KeyVal=[VectSerialMode,csv]
MMVII EditProfile KeyVal=[PdfOpen,evince]
MMVII EditProfile KeyVal=[Empty,,]
\end{verbatim}

The command writes the modified profile file after applying the change.  Values are
stored as strings in XML and converted later by the code that reads the key.

The optional parameter \FileSty{DelKey} removes a key from the selected profile:

\begin{verbatim}
MMVII EditProfile DelKey=PdfOpen
\end{verbatim}

Removing a key from a user profile does not remove it from the distributed default
profile.  At the next synchronization, keys present in
\FileSty{MMVII-Default-Profile.xml} are restored with their default values unless
the user profile overrides them.


 %=================================================================================

\section{Command options}

\begin{center}
\begin{tabular}{|p{0.20\linewidth}|p{0.18\linewidth}|p{0.52\linewidth}|}
\hline
Option & Type & Meaning \\
\hline
\FileSty{SetCurrent} & profile name & Name of the profile to make current \\
\hline
\FileSty{KeyVal} & vector of two strings & Set a profile value as \FileSty{[Key,Value]} \\
\hline
\FileSty{DelKey} & profile key & Delete one key from the selected profile \\
\hline
\end{tabular}
\end{center}

If none of these options is used, \TheCmd\ only displays the current profile
information and does not modify any profile file.


 %=================================================================================

\section{Using a profile for one command}

Every MMVII command has a global \FileSty{Profile} option.  It selects a profile
only for the current execution:

\begin{verbatim}
MMVII SomeCommand ... Profile=Test
\end{verbatim}

This does not modify \FileSty{MMVII-Current-Profile.xml}.  It is useful for scripts
or for testing profile-specific behavior without changing the user's default
profile.

Command-line options remain authoritative for one command execution.  For example,
\FileSty{NbProcMax} is only applied when the command did not explicitly initialize
the global \FileSty{NbProc} option.  When both are present, the explicit command-line
option keeps priority.


 %=================================================================================

\section{User profile mechanism}

MMVII user profiles store per-user defaults outside the project data directory.
They are loaded by every MMVII command.  A command normally uses the current
profile, but it can also select another profile for one execution with the global
\FileSty{Profile} option.

\subsection{Configuration directory}

User profiles are stored in an operating-system-specific MMVII configuration
directory:

\begin{center}
\begin{tabular}{|l|p{0.72\linewidth}|}
\hline
System & Directory \\
\hline
Linux & \FileSty{\$XDG\_CONFIG\_HOME/MMVII} when \FileSty{XDG\_CONFIG\_HOME} is an absolute path, otherwise \FileSty{\$HOME/.config/MMVII} \\
\hline
macOS & \FileSty{\$HOME/Library/Application Support/MMVII} \\
\hline
Windows & \FileSty{\%APPDATA\%/MMVII} \\
\hline
\end{tabular}
\end{center}

MMVII creates this directory when it needs to write the current profile selection
or a user profile file.

\subsection{Profile files}

The profile mechanism uses three file families:

\begin{table}[h]
\centering
\caption{Configuration files and their locations}
\begin{tabular}{|p{0.42\linewidth}|p{0.48\linewidth}|}
\hline
File & Location \\
\hline
{\tt MMVII-Default-Profile.xml} &
{\tt MMVII-LocalParameters/} installation directory \\
\hline
{\tt MMVII-Current-Profile.xml} &
user configuration directory \\
\hline
{\tt MMVII-profile-\textless Name\textgreater.xml} &
user configuration directory \\
\hline
\end{tabular}
\end{table}

\begin{table}[h]
\centering
\caption{Configuration files and their roles}
\begin{tabular}{|p{0.42\linewidth}|p{0.48\linewidth}|}
\hline
File & Role \\
\hline
{\tt MMVII-Default-Profile.xml} &
Default key/value set distributed with MMVII \\
\hline
{\tt MMVII-Current-Profile.xml} &
Stores the name of the current profile. \\
\hline
{\tt MMVII-profile-\textless Name\textgreater.xml} &
Stores the user-defined values for profile {\tt \textless Name\textgreater}. \\
\hline
\end{tabular}
\end{table}


If \FileSty{MMVII-Current-Profile.xml} does not exist, MMVII uses the profile name
\FileSty{Default}.  When a profile is loaded, MMVII first reads
\FileSty{MMVII-Default-Profile.xml}, then overlays values from
\FileSty{MMVII-profile-<Name>.xml} if that file already exists.  It writes the
merged result back to the user profile file.  This keeps user profile files
synchronized when new default keys are added by MMVII.

\subsection{Standard keys}

The default profile currently defines these keys:

\begin{center}
\begin{tabular}{|p{0.22\linewidth}|p{0.14\linewidth}|p{0.54\linewidth}|}
\hline
Key & Default & Meaning \\
\hline
\FileSty{NbProcMax} & \FileSty{4} & Upper limit for the number of parallel processes used by commands when \FileSty{NbProc} was not explicitly provided \\
\hline
\FileSty{TaggedSerialMode} & \FileSty{xml} & Default serialization mode for tagged objects \\
\hline
\FileSty{VectSerialMode} & \FileSty{csv} & Default serialization mode for vector-like serialized data \\
\hline
\FileSty{PdfOpen} & empty string & Command used by MMVII to open generated PDF files, when configured \\
\hline
\end{tabular}
\end{center}

If a key is absent, MMVII uses the caller-provided default value and emits a user
warning.  The profile is initialized early enough for help and argument
specification generation, so profile names and profile keys can be exposed as
completion or GUI metadata through argument types tagged as
\FileSty{eTA2007::Profile} and \FileSty{eTA2007::ProfileKey}.

\end{document}

